% --- Archon physics formalization source begin ---
% archon:physics
% archon:covers PhyXMiniProblems/problem_phyx_mini_0740.lean
% archon:source-report reports/phyx_mini/problem_phyx_mini_0740.source.json

\chapter{Physics problem phyx\_mini\_0740}
\label{ch:PhyXMiniProblems_problem_phyx_mini_0740}

\paragraph{Problem source.}
\#\# Physical scenario

You throw a ball toward a wall at speed \$25.0\textbackslash{} m/s\$ and at angle \$\textbackslash{}theta\_0 = 40.0\textasciicircum{}\textbackslash{}circ\$ above the horizontal. The wall is distance \$d = 22.0\textbackslash{} m\$ from the release point of the ball.

\#\# Question

How far above the release point does the ball hit the wall?

\#\# Answer choices

- A: A: 8.0 m

- B: B: 10.0 m

- C: C: 12.0 m

- D: D: 14.0 m

\#\# Figure caption (auxiliary; use the image as primary evidence)

The image is a physics illustration showing projectile motion. It features a person on the left who appears to be throwing an object at an angle \textbackslash{}(\textbackslash{}theta\_0\textbackslash{}) above the horizontal. The trajectory of the object is shown as a curved, dotted path, going upwards and then curving downwards towards the right.

On the right side, there is a vertical brick wall. The horizontal distance from the person to the wall is marked as \textbackslash{}(d\textbackslash{}). The path of the object suggests it is aimed to potentially reach or hit the wall.

The surface below the person, the trajectory, and the wall is textured to resemble some form of ground pattern. The diagram emphasizes the initial angle of projection and the horizontal distance to the target.

\#\# Dataset metadata

Kinematics; Spatial Relation Reasoning, Multi-Formula Reasoning

\paragraph{Recorded answer/context.}
C: C: 12.0 m

\paragraph{Figure/image path.}
/root/proposal\_for\_physic/science-mango/phyx\_mini\_run/phyx\_data/test\_image/740.png

\paragraph{Formalization target.}
create a compiling Lean file with sorry bodies at `PhyXMiniProblems/problem\_phyx\_mini\_0740.lean`. The Lean declarations must preserve the physical quantities, dimensions or dimensional roles, figure labels, governing-law hypotheses, and final relation expressed by this problem.
Use Mathlib/Physlib names found through LeanExplore where available. If a physics API is missing, introduce faithful local abstractions rather than scalar placeholder aliases.

\begin{theorem}[Physics formalization target]
\label{thm:physics:phyx_mini_0740:target}
\lean{PhyXMiniProblems.ProblemPhyXMini0740.problem_phyx_mini_0740}
\uses{def:physics:phyx-mini-0740:phyxminiproblems-problemphyxmini0740-lengthinmeters, def:physics:phyx-mini-0740:phyxminiproblems-problemphyxmini0740-ballwallprojectilesetup, def:physics:phyx-mini-0740:phyxminiproblems-problemphyxmini0740-matchesprimaryfigure, def:physics:phyx-mini-0740:phyxminiproblems-problemphyxmini0740-matchesproblemdescription, def:physics:phyx-mini-0740:phyxminiproblems-problemphyxmini0740-usesstandardnearearthgravity, def:physics:phyx-mini-0740:phyxminiproblems-problemphyxmini0740-hasphysicalprojectileparameters, def:physics:phyx-mini-0740:phyxminiproblems-problemphyxmini0740-satisfiesidealprojectilekinematics, def:physics:phyx-mini-0740:phyxminiproblems-problemphyxmini0740-describeswallimpact, def:physics:phyx-mini-0740:phyxminiproblems-problemphyxmini0740-computedimpactheightinmeters, def:physics:phyx-mini-0740:phyxminiproblems-problemphyxmini0740-matchesdisplayedheight, def:physics:phyx-mini-0740:phyxminiproblems-problemphyxmini0740-isuniqueclosestdisplayedheight}
The assigned autoformalize agent should translate this physics problem into Lean declarations in the covered file, with theorem and lemma proof bodies written as `by sorry`.
\end{theorem}
\begin{proof}
This is an autoformalization task, not a proof task. Produce statements that can later be proved by the physics prover without weakening the physical model.
\end{proof}

% --- Archon declaration topology begin ---
\section{Declaration topology}
\begin{definition}[acceleration Dimension]
  \label{def:physics:phyx-mini-0740:phyxminiproblems-problemphyxmini0740-accelerationdimension}
  \lean{PhyXMiniProblems.ProblemPhyXMini0740.accelerationDimension}
  Acceleration has physical dimension L T\textasciicircum{}-2
\end{definition}
\begin{proof}
  This is the declared model definition of acceleration Dimension.
\end{proof}
\begin{definition}[Length Quantity]
  \label{def:physics:phyx-mini-0740:phyxminiproblems-problemphyxmini0740-lengthquantity}
  \lean{PhyXMiniProblems.ProblemPhyXMini0740.LengthQuantity}
  A nonnegative, unit-independent physical length
\end{definition}
\begin{proof}
  This is the declared model definition of Length Quantity.
\end{proof}
\begin{definition}[Signed Length Quantity]
  \label{def:physics:phyx-mini-0740:phyxminiproblems-problemphyxmini0740-signedlengthquantity}
  \lean{PhyXMiniProblems.ProblemPhyXMini0740.SignedLengthQuantity}
  A signed physical displacement along one Cartesian axis
\end{definition}
\begin{proof}
  This is the declared model definition of Signed Length Quantity.
\end{proof}
\begin{definition}[Time Quantity]
  \label{def:physics:phyx-mini-0740:phyxminiproblems-problemphyxmini0740-timequantity}
  \lean{PhyXMiniProblems.ProblemPhyXMini0740.TimeQuantity}
  A nonnegative, unit-independent elapsed duration
\end{definition}
\begin{proof}
  This is the declared model definition of Time Quantity.
\end{proof}
\begin{definition}[Speed Quantity]
  \label{def:physics:phyx-mini-0740:phyxminiproblems-problemphyxmini0740-speedquantity}
  \lean{PhyXMiniProblems.ProblemPhyXMini0740.SpeedQuantity}
  A nonnegative, unit-independent speed magnitude
\end{definition}
\begin{proof}
  This is the declared model definition of Speed Quantity.
\end{proof}
\begin{definition}[Acceleration Quantity]
  \label{def:physics:phyx-mini-0740:phyxminiproblems-problemphyxmini0740-accelerationquantity}
  \lean{PhyXMiniProblems.ProblemPhyXMini0740.AccelerationQuantity}
  \uses{def:physics:phyx-mini-0740:phyxminiproblems-problemphyxmini0740-accelerationdimension}
  A nonnegative magnitude of gravitational acceleration
\end{definition}
\begin{proof}
  This is the declared model definition of Acceleration Quantity.
\end{proof}
\begin{definition}[length In Meters]
  \label{def:physics:phyx-mini-0740:phyxminiproblems-problemphyxmini0740-lengthinmeters}
  \lean{PhyXMiniProblems.ProblemPhyXMini0740.lengthInMeters}
  \uses{def:physics:phyx-mini-0740:phyxminiproblems-problemphyxmini0740-lengthquantity}
  Read a nonnegative physical length in metres
\end{definition}
\begin{proof}
  This is the declared model definition of length In Meters.
\end{proof}
\begin{definition}[signed Length In Meters]
  \label{def:physics:phyx-mini-0740:phyxminiproblems-problemphyxmini0740-signedlengthinmeters}
  \lean{PhyXMiniProblems.ProblemPhyXMini0740.signedLengthInMeters}
  \uses{def:physics:phyx-mini-0740:phyxminiproblems-problemphyxmini0740-signedlengthquantity}
  Read a signed displacement in metres
\end{definition}
\begin{proof}
  This is the declared model definition of signed Length In Meters.
\end{proof}
\begin{definition}[time In Seconds]
  \label{def:physics:phyx-mini-0740:phyxminiproblems-problemphyxmini0740-timeinseconds}
  \lean{PhyXMiniProblems.ProblemPhyXMini0740.timeInSeconds}
  \uses{def:physics:phyx-mini-0740:phyxminiproblems-problemphyxmini0740-timequantity}
  Read a physical duration in seconds
\end{definition}
\begin{proof}
  This is the declared model definition of time In Seconds.
\end{proof}
\begin{definition}[speed In Meters Per Second]
  \label{def:physics:phyx-mini-0740:phyxminiproblems-problemphyxmini0740-speedinmeterspersecond}
  \lean{PhyXMiniProblems.ProblemPhyXMini0740.speedInMetersPerSecond}
  \uses{def:physics:phyx-mini-0740:phyxminiproblems-problemphyxmini0740-speedquantity}
  Read a physical speed magnitude in metres per second
\end{definition}
\begin{proof}
  This is the declared model definition of speed In Meters Per Second.
\end{proof}
\begin{definition}[acceleration In Meters Per Second Squared]
  \label{def:physics:phyx-mini-0740:phyxminiproblems-problemphyxmini0740-accelerationinmeterspersecondsquared}
  \lean{PhyXMiniProblems.ProblemPhyXMini0740.accelerationInMetersPerSecondSquared}
  \uses{def:physics:phyx-mini-0740:phyxminiproblems-problemphyxmini0740-accelerationquantity}
  Read an acceleration magnitude in metres per second squared
\end{definition}
\begin{proof}
  This is the declared model definition of acceleration In Meters Per Second Squared.
\end{proof}
\begin{definition}[Projectile Model]
  \label{def:physics:phyx-mini-0740:phyxminiproblems-problemphyxmini0740-projectilemodel}
  \lean{PhyXMiniProblems.ProblemPhyXMini0740.ProjectileModel}
  The idealization used for the flight between release and wall impact
\end{definition}
\begin{proof}
  This is the declared model definition of Projectile Model.
\end{proof}
\begin{definition}[Figure Object]
  \label{def:physics:phyx-mini-0740:phyxminiproblems-problemphyxmini0740-figureobject}
  \lean{PhyXMiniProblems.ProblemPhyXMini0740.FigureObject}
  Literal objects visibly represented in image 740.png
\end{definition}
\begin{proof}
  This is the declared model definition of Figure Object.
\end{proof}
\begin{definition}[Figure Label]
  \label{def:physics:phyx-mini-0740:phyxminiproblems-problemphyxmini0740-figurelabel}
  \lean{PhyXMiniProblems.ProblemPhyXMini0740.FigureLabel}
  Mathematical labels visibly printed in the primary image
\end{definition}
\begin{proof}
  This is the declared model definition of Figure Label.
\end{proof}
\begin{definition}[Figure Anchor]
  \label{def:physics:phyx-mini-0740:phyxminiproblems-problemphyxmini0740-figureanchor}
  \lean{PhyXMiniProblems.ProblemPhyXMini0740.FigureAnchor}
  Schematic anchors used to retain the spatial relations in the raster
\end{definition}
\begin{proof}
  This is the declared model definition of Figure Anchor.
\end{proof}
\begin{definition}[Ball Wall Figure]
  \label{def:physics:phyx-mini-0740:phyxminiproblems-problemphyxmini0740-ballwallfigure}
  \lean{PhyXMiniProblems.ProblemPhyXMini0740.BallWallFigure}
  \uses{def:physics:phyx-mini-0740:phyxminiproblems-problemphyxmini0740-lengthquantity, def:physics:phyx-mini-0740:phyxminiproblems-problemphyxmini0740-figureobject, def:physics:phyx-mini-0740:phyxminiproblems-problemphyxmini0740-figurelabel, def:physics:phyx-mini-0740:phyxminiproblems-problemphyxmini0740-figureanchor}
  The primary-image transcription. Its real coordinates are dimensionless drawing coordinates rather than physical measurements. The two labelled physical quantities remain dimensionful and are connected to the experiment by MatchesPrimaryFigure
\end{definition}
\begin{proof}
  This is the declared model definition of Ball Wall Figure.
\end{proof}
\begin{definition}[Ball Wall Projectile Setup]
  \label{def:physics:phyx-mini-0740:phyxminiproblems-problemphyxmini0740-ballwallprojectilesetup}
  \lean{PhyXMiniProblems.ProblemPhyXMini0740.BallWallProjectileSetup}
  \uses{def:physics:phyx-mini-0740:phyxminiproblems-problemphyxmini0740-lengthquantity, def:physics:phyx-mini-0740:phyxminiproblems-problemphyxmini0740-signedlengthquantity, def:physics:phyx-mini-0740:phyxminiproblems-problemphyxmini0740-timequantity, def:physics:phyx-mini-0740:phyxminiproblems-problemphyxmini0740-speedquantity, def:physics:phyx-mini-0740:phyxminiproblems-problemphyxmini0740-accelerationquantity, def:physics:phyx-mini-0740:phyxminiproblems-problemphyxmini0740-projectilemodel, def:physics:phyx-mini-0740:phyxminiproblems-problemphyxmini0740-ballwallfigure}
  Independent physical observables of the experiment. The coordinate functions give displacement from the release point as a function of physical elapsed time. In particular, impactHeightAboveRelease and impactTime are not defined from the requested answer or from an eliminated projectile formula
\end{definition}
\begin{proof}
  This is the declared model definition of Ball Wall Projectile Setup.
\end{proof}
\begin{definition}[Matches Primary Figure]
  \label{def:physics:phyx-mini-0740:phyxminiproblems-problemphyxmini0740-matchesprimaryfigure}
  \lean{PhyXMiniProblems.ProblemPhyXMini0740.MatchesPrimaryFigure}
  \uses{def:physics:phyx-mini-0740:phyxminiproblems-problemphyxmini0740-ballwallprojectilesetup}
  The raster shows the release point to the left of the vertical wall, a ball above and to the right of release but still left of the wall, the horizontal arrow labelled d, and the angle theta\_0 above a horizontal reference. The image itself does not display a numerical impact height
\end{definition}
\begin{proof}
  This is the declared model definition of Matches Primary Figure.
\end{proof}
\begin{definition}[degrees]
  \label{def:physics:phyx-mini-0740:phyxminiproblems-problemphyxmini0740-degrees}
  \lean{PhyXMiniProblems.ProblemPhyXMini0740.degrees}
  Convert a degree readout to a physical angle modulo 2 * pi
\end{definition}
\begin{proof}
  This is the declared model definition of degrees.
\end{proof}
\begin{definition}[Matches Problem Description]
  \label{def:physics:phyx-mini-0740:phyxminiproblems-problemphyxmini0740-matchesproblemdescription}
  \lean{PhyXMiniProblems.ProblemPhyXMini0740.MatchesProblemDescription}
  \uses{def:physics:phyx-mini-0740:phyxminiproblems-problemphyxmini0740-lengthinmeters, def:physics:phyx-mini-0740:phyxminiproblems-problemphyxmini0740-speedinmeterspersecond, def:physics:phyx-mini-0740:phyxminiproblems-problemphyxmini0740-ballwallprojectilesetup, def:physics:phyx-mini-0740:phyxminiproblems-problemphyxmini0740-degrees}
  Numerical and qualitative data stated in the problem prose. No field here mentions the impact height or an answer choice
\end{definition}
\begin{proof}
  This is the declared model definition of Matches Problem Description.
\end{proof}
\begin{definition}[Uses Standard Near Earth Gravity]
  \label{def:physics:phyx-mini-0740:phyxminiproblems-problemphyxmini0740-usesstandardnearearthgravity}
  \lean{PhyXMiniProblems.ProblemPhyXMini0740.UsesStandardNearEarthGravity}
  \uses{def:physics:phyx-mini-0740:phyxminiproblems-problemphyxmini0740-accelerationinmeterspersecondsquared, def:physics:phyx-mini-0740:phyxminiproblems-problemphyxmini0740-ballwallprojectilesetup}
  The standard near-Earth gravitational calibration implicit in the textbook answer. It is independent of the requested impact height
\end{definition}
\begin{proof}
  This is the declared model definition of Uses Standard Near Earth Gravity.
\end{proof}
\begin{definition}[Has Physical Projectile Parameters]
  \label{def:physics:phyx-mini-0740:phyxminiproblems-problemphyxmini0740-hasphysicalprojectileparameters}
  \lean{PhyXMiniProblems.ProblemPhyXMini0740.HasPhysicalProjectileParameters}
  \uses{def:physics:phyx-mini-0740:phyxminiproblems-problemphyxmini0740-lengthinmeters, def:physics:phyx-mini-0740:phyxminiproblems-problemphyxmini0740-timeinseconds, def:physics:phyx-mini-0740:phyxminiproblems-problemphyxmini0740-speedinmeterspersecond, def:physics:phyx-mini-0740:phyxminiproblems-problemphyxmini0740-accelerationinmeterspersecondsquared, def:physics:phyx-mini-0740:phyxminiproblems-problemphyxmini0740-ballwallprojectilesetup}
  Strict positivity and the up-and-right launch branch shown by the figure. These assumptions support cancellation of horizontal speed and exclude a spurious backwards or zero-time incidence
\end{definition}
\begin{proof}
  This is the declared model definition of Has Physical Projectile Parameters.
\end{proof}
\begin{definition}[Satisfies Ideal Projectile Kinematics]
  \label{def:physics:phyx-mini-0740:phyxminiproblems-problemphyxmini0740-satisfiesidealprojectilekinematics}
  \lean{PhyXMiniProblems.ProblemPhyXMini0740.SatisfiesIdealProjectileKinematics}
  \uses{def:physics:phyx-mini-0740:phyxminiproblems-problemphyxmini0740-timequantity, def:physics:phyx-mini-0740:phyxminiproblems-problemphyxmini0740-signedlengthinmeters, def:physics:phyx-mini-0740:phyxminiproblems-problemphyxmini0740-timeinseconds, def:physics:phyx-mini-0740:phyxminiproblems-problemphyxmini0740-speedinmeterspersecond, def:physics:phyx-mini-0740:phyxminiproblems-problemphyxmini0740-accelerationinmeterspersecondsquared, def:physics:phyx-mini-0740:phyxminiproblems-problemphyxmini0740-ballwallprojectilesetup}
  Ideal no-drag projectile motion in coherent SI component readouts. Horizontal velocity is constant; vertical acceleration is the uniform downward g. The two equations are governing laws for every elapsed duration and contain no wall-impact height formula or answer-choice value
\end{definition}
\begin{proof}
  This is the declared model definition of Satisfies Ideal Projectile Kinematics.
\end{proof}
\begin{definition}[Describes Wall Impact]
  \label{def:physics:phyx-mini-0740:phyxminiproblems-problemphyxmini0740-describeswallimpact}
  \lean{PhyXMiniProblems.ProblemPhyXMini0740.DescribesWallImpact}
  \uses{def:physics:phyx-mini-0740:phyxminiproblems-problemphyxmini0740-lengthinmeters, def:physics:phyx-mini-0740:phyxminiproblems-problemphyxmini0740-signedlengthinmeters, def:physics:phyx-mini-0740:phyxminiproblems-problemphyxmini0740-ballwallprojectilesetup}
  Incidence and measurement relations for the independently stored wall-impact event. They only say that the impact lies at the wall and that the named height is its vertical displacement above release
\end{definition}
\begin{proof}
  This is the declared model definition of Describes Wall Impact.
\end{proof}
\begin{definition}[predicted Impact Time In Seconds]
  \label{def:physics:phyx-mini-0740:phyxminiproblems-problemphyxmini0740-predictedimpacttimeinseconds}
  \lean{PhyXMiniProblems.ProblemPhyXMini0740.predictedImpactTimeInSeconds}
  \uses{def:physics:phyx-mini-0740:phyxminiproblems-problemphyxmini0740-lengthinmeters, def:physics:phyx-mini-0740:phyxminiproblems-problemphyxmini0740-speedinmeterspersecond, def:physics:phyx-mini-0740:phyxminiproblems-problemphyxmini0740-ballwallprojectilesetup}
  The wall-arrival time predicted by uniform horizontal motion
\end{definition}
\begin{proof}
  This is the declared model definition of predicted Impact Time In Seconds.
\end{proof}
\begin{definition}[predicted Impact Height In Meters]
  \label{def:physics:phyx-mini-0740:phyxminiproblems-problemphyxmini0740-predictedimpactheightinmeters}
  \lean{PhyXMiniProblems.ProblemPhyXMini0740.predictedImpactHeightInMeters}
  \uses{def:physics:phyx-mini-0740:phyxminiproblems-problemphyxmini0740-speedinmeterspersecond, def:physics:phyx-mini-0740:phyxminiproblems-problemphyxmini0740-accelerationinmeterspersecondsquared, def:physics:phyx-mini-0740:phyxminiproblems-problemphyxmini0740-ballwallprojectilesetup, def:physics:phyx-mini-0740:phyxminiproblems-problemphyxmini0740-predictedimpacttimeinseconds}
  The height above release predicted after substituting the derived wall-arrival time into the vertical projectile equation. This expression does not define the independent physical impact-height field
\end{definition}
\begin{proof}
  This is the declared model definition of predicted Impact Height In Meters.
\end{proof}
\begin{definition}[computed Impact Height In Meters]
  \label{def:physics:phyx-mini-0740:phyxminiproblems-problemphyxmini0740-computedimpactheightinmeters}
  \lean{PhyXMiniProblems.ProblemPhyXMini0740.computedImpactHeightInMeters}
  \uses{def:physics:phyx-mini-0740:phyxminiproblems-problemphyxmini0740-degrees}
  The exact real expression obtained from the displayed 25, 40 degrees, 22, and the conventional 9.8 calibration. Its numerical value is about 12.0, but it is not definitionally set equal to 12
\end{definition}
\begin{proof}
  This is the declared model definition of computed Impact Height In Meters.
\end{proof}
\begin{definition}[Answer Choice]
  \label{def:physics:phyx-mini-0740:phyxminiproblems-problemphyxmini0740-answerchoice}
  \lean{PhyXMiniProblems.ProblemPhyXMini0740.AnswerChoice}
  The four answer labels printed with the problem
\end{definition}
\begin{proof}
  This is the declared model definition of Answer Choice.
\end{proof}
\begin{definition}[displayed Height In Meters]
  \label{def:physics:phyx-mini-0740:phyxminiproblems-problemphyxmini0740-answerchoice-displayedheightinmeters}
  \lean{PhyXMiniProblems.ProblemPhyXMini0740.AnswerChoice.displayedHeightInMeters}
  \uses{def:physics:phyx-mini-0740:phyxminiproblems-problemphyxmini0740-answerchoice}
  Height in metres printed beside each displayed answer label
\end{definition}
\begin{proof}
  This is the declared model definition of displayed Height In Meters.
\end{proof}
\begin{definition}[recorded Dataset Answer]
  \label{def:physics:phyx-mini-0740:phyxminiproblems-problemphyxmini0740-recordeddatasetanswer}
  \lean{PhyXMiniProblems.ProblemPhyXMini0740.recordedDatasetAnswer}
  \uses{def:physics:phyx-mini-0740:phyxminiproblems-problemphyxmini0740-answerchoice}
  The answer label recorded in the supplied dataset metadata
\end{definition}
\begin{proof}
  This is the declared model definition of recorded Dataset Answer.
\end{proof}
\begin{definition}[Matches Displayed Height]
  \label{def:physics:phyx-mini-0740:phyxminiproblems-problemphyxmini0740-matchesdisplayedheight}
  \lean{PhyXMiniProblems.ProblemPhyXMini0740.MatchesDisplayedHeight}
  \uses{def:physics:phyx-mini-0740:phyxminiproblems-problemphyxmini0740-lengthquantity, def:physics:phyx-mini-0740:phyxminiproblems-problemphyxmini0740-lengthinmeters, def:physics:phyx-mini-0740:phyxminiproblems-problemphyxmini0740-answerchoice}
  Agreement with a displayed whole-metre height to within half a metre
\end{definition}
\begin{proof}
  This is the declared model definition of Matches Displayed Height.
\end{proof}
\begin{definition}[Is Unique Closest Displayed Height]
  \label{def:physics:phyx-mini-0740:phyxminiproblems-problemphyxmini0740-isuniqueclosestdisplayedheight}
  \lean{PhyXMiniProblems.ProblemPhyXMini0740.IsUniqueClosestDisplayedHeight}
  \uses{def:physics:phyx-mini-0740:phyxminiproblems-problemphyxmini0740-lengthquantity, def:physics:phyx-mini-0740:phyxminiproblems-problemphyxmini0740-lengthinmeters, def:physics:phyx-mini-0740:phyxminiproblems-problemphyxmini0740-answerchoice}
  A choice is strictly closer to the physical height than every alternative
\end{definition}
\begin{proof}
  This is the declared model definition of Is Unique Closest Displayed Height.
\end{proof}
\begin{lemma}[impact Time eq predicted]
  \label{lem:physics:phyx-mini-0740:phyxminiproblems-problemphyxmini0740-impacttime-eq-predicted}
  \lean{PhyXMiniProblems.ProblemPhyXMini0740.impactTime_eq_predicted}
  \uses{def:physics:phyx-mini-0740:phyxminiproblems-problemphyxmini0740-timeinseconds, def:physics:phyx-mini-0740:phyxminiproblems-problemphyxmini0740-ballwallprojectilesetup, def:physics:phyx-mini-0740:phyxminiproblems-problemphyxmini0740-hasphysicalprojectileparameters, def:physics:phyx-mini-0740:phyxminiproblems-problemphyxmini0740-satisfiesidealprojectilekinematics, def:physics:phyx-mini-0740:phyxminiproblems-problemphyxmini0740-describeswallimpact, def:physics:phyx-mini-0740:phyxminiproblems-problemphyxmini0740-predictedimpacttimeinseconds}
  The wall-incidence condition and horizontal kinematics determine the positive impact time. This is a derived result, not part of either law structure
\end{lemma}
\begin{proof}
  The conclusion follows from the governing relations and figure readouts named in the statement of impact Time eq predicted.
\end{proof}
\begin{lemma}[impact Height eq predicted]
  \label{lem:physics:phyx-mini-0740:phyxminiproblems-problemphyxmini0740-impactheight-eq-predicted}
  \lean{PhyXMiniProblems.ProblemPhyXMini0740.impactHeight_eq_predicted}
  \uses{def:physics:phyx-mini-0740:phyxminiproblems-problemphyxmini0740-lengthinmeters, def:physics:phyx-mini-0740:phyxminiproblems-problemphyxmini0740-ballwallprojectilesetup, def:physics:phyx-mini-0740:phyxminiproblems-problemphyxmini0740-hasphysicalprojectileparameters, def:physics:phyx-mini-0740:phyxminiproblems-problemphyxmini0740-satisfiesidealprojectilekinematics, def:physics:phyx-mini-0740:phyxminiproblems-problemphyxmini0740-describeswallimpact, def:physics:phyx-mini-0740:phyxminiproblems-problemphyxmini0740-predictedimpactheightinmeters}
  Substitution into the vertical kinematics gives the general impact-height prediction without using any displayed answer
\end{lemma}
\begin{proof}
  The conclusion follows from the governing relations and figure readouts named in the statement of impact Height eq predicted.
\end{proof}
% --- Archon declaration topology end ---
% --- Archon physics formalization source end ---
